\begin{proof}
Let \(x\) and \(y\) be vectors in an inner–product space \((V,\langle\cdot,\cdot\rangle)\).
The norm induced by the inner product is  

\[
\|v\|=\sqrt{\langle v,v\rangle}\qquad (v\in V).
\tag{1}
\]

\medskip
\noindent\textbf{1.  Reduction to the squared inequality.}
Both \(\|x+y\|\) and \(\|x\|+\|y\|\) are non‑negative, hence the triangle inequality  

\[
\|x+y\|\le \|x\|+\|y\|
\tag{2}
\]

is equivalent to  

\[
\|x+y\|^{2}\le (\|x\|+\|y\|)^{2}.
\tag{3}
\]

\medskip
\noindent\textbf{2.  Expansion of \(\|x+y\|^{2}\).}
Using bilinearity of the inner product and \(\langle y,x\rangle
=\overline{\langle x,y\rangle}\),

\[
\begin{aligned}
\|x+y\|^{2}
   &=\langle x+y,\,x+y\rangle \\[2pt]
   &=\langle x,x\rangle+\langle x,y\rangle
     +\langle y,x\rangle+\langle y,y\rangle \\[2pt]
   &=\|x\|^{2}+2\operatorname{Re}\langle x,y\rangle+\|y\|^{2}.
\end{aligned}
\tag{4}
\]

\medskip
\noindent\textbf{3.  Expansion of the right–hand side of (3).}

\[
(\|x\|+\|y\|)^{2}= \|x\|^{2}+2\|x\|\|y\|+\|y\|^{2}.
\tag{5}
\]

\medskip
\noindent\textbf{4.  Application of Cauchy–Schwarz.}
The Cauchy–Schwarz inequality gives  

\[
|\langle x,y\rangle|\le \|x\|\,\|y\|.
\tag{6}
\]

Since \(\operatorname{Re}\langle x,y\rangle\le |\langle x,y\rangle|\), we obtain  

\[
2\operatorname{Re}\langle x,y\rangle\le 2\|x\|\,\|y\|.
\tag{7}
\]

\medskip
\noindent\textbf{5.  Comparison of (4) and (5).}
Insert the estimate (7) into (4):

\[
\begin{aligned}
\|x+y\|^{2}
   &=\|x\|^{2}+2\operatorname{Re}\langle x,y\rangle+\|y\|^{2}\\
   &\le \|x\|^{2}+2\|x\|\,\|y\|+\|y\|^{2}\\
   &= (\|x\|+\|y\|)^{2}.
\end{aligned}
\]

Thus (3) holds.

\medskip
\noindent\textbf{6.  Return to the original inequality.}
Both sides of (3) are non‑negative; taking square roots yields (2):

\[
\boxed{\;\|x+y\|\le \|x\|+\|y\|\;}.
\]

\medskip
\noindent\textbf{7.  Equality condition.}
Equality in the chain above occurs precisely when equality holds in both
\(\operatorname{Re}\langle x,y\rangle\le |\langle x,y\rangle|\) and in the
Cauchy–Schwarz inequality (6).

\begin{itemize}
\item \(\operatorname{Re}\langle x,y\rangle = |\langle x,y\rangle|\) means that
\(\langle x,y\rangle\) is a non‑negative real number.
\item Equality in Cauchy–Schwarz holds iff \(x\) and \(y\) are linearly dependent,
i.e. \(y=\lambda x\) for some scalar \(\lambda\in\mathbb{C}\).
\end{itemize}
Combining the two requirements forces \(\lambda\) to be a \emph{real non‑negative}
scalar. Hence, for non‑zero vectors,

\[
\|x+y\| = \|x\|+\|y\|\quad\Longleftrightarrow\quad y=\lambda x\ \text{with}\ \lambda\ge 0.
\]

The statement also holds trivially when \(x=0\) or \(y=0\).

Consequently, the triangle inequality is proved, and equality occurs exactly
when the vectors are collinear and point in the same (or zero) direction. \(\square\)
\end{proof}